\documentclass[11pt]{article}
\usepackage{../../../tex/lecturestyle}

\title{Kth Term of a Linearly Recurrent Sequence}
\author{}
\date{}

\begin{document}
\thispagestyle{fancy}

\lecturetitleband{Polynomials \& Formal Power Series}{Kth Term of a Linearly Recurrent Sequence}{6}

\section{Overview}
Given a degree-$d$ constant-coefficient linear recurrence
$A_i = c_1 A_{i-1} + \cdots + c_d A_{i-d}$ ($i \ge d$) together with initial
values $A_0, \ldots, A_{d-1}$ and a target index $K$, how quickly can
$A_K$ be found? Matrix exponentiation answers this in
$O(d^3 \log K)$. This lecture covers two classical algorithms that do
better, $O(M(d) \log K)$, where $M(d)$ is the cost of multiplying two
degree-$d$ polynomials ($O(d^2)$ naively, $O(d \log d)$ with FFT-based
multiplication over a suitable ring): \textbf{Fiduccia's algorithm}, which
reduces the problem to computing $x^K$ modulo a fixed polynomial, and the
\textbf{Bostan--Mori algorithm}, which works directly with the sequence's
rational generating function. Assumes the companion lecture on linearly
recurrent sequences (generating functions as $P(x)/Q(x)$) and basic
polynomial multiplication/remainder.

Throughout, arithmetic is over a fixed commutative ring with unit
(think $\mathbb{Z}/m\mathbb{Z}$ or $\mathbb{F}_p$), with ring operations
taken as $O(1)$.

\section{Fiduccia's Algorithm}

\subsection{Recurrences as polynomial remainders}
Unrolling the recurrence expresses $A_i$ as a linear combination of
$A_0, \ldots, A_{d-1}$: substitute $A_{i} = c_1 A_{i-1} + \cdots$
repeatedly until every index drops below $d$. This is formally identical
to reducing $x^i$ modulo a fixed polynomial.

\begin{definitionbox}[(Characteristic polynomial)]
For the recurrence above, its characteristic polynomial is
\[
  \Gamma(x) = x^d - c_1 x^{d-1} - c_2 x^{d-2} - \cdots - c_d .
\]
\end{definitionbox}

\begin{theorembox}
Let $A$ satisfy the recurrence with characteristic polynomial $\Gamma(x)$.
If $x^i \bmod \Gamma(x) = r_0 + r_1 x + \cdots + r_{d-1}x^{d-1}$, then
\[
  A_i = r_0 A_0 + r_1 A_1 + \cdots + r_{d-1} A_{d-1}.
\]

\emph{Why.} Both ``replace $A_i$ using the recurrence'' and ``replace
$x^d$ using $x^d \equiv c_1 x^{d-1} + \cdots + c_d \pmod{\Gamma(x)}$''
apply the exact same substitution rule with the same coefficients, so
tracking one tracks the other. \qed
\end{theorembox}

\begin{examplebox}
For $A_i = A_{i-1}+A_{i-2}+A_{i-3}$ with $A_0=3, A_1=2, A_2=1$: unrolling
gives $A_3=6$, $A_4=9$, $A_5=16$; matching this against
$x^5 \bmod (x^3-x^2-x-1) = 4x^2+3x+2$ gives
$A_5 = 4A_2+3A_1+2A_0 = 4{\cdot}1+3{\cdot}2+2{\cdot}3 = 16$, confirming the
theorem.
\end{examplebox}

\subsection{The algorithm}
$x^K \bmod \Gamma(x)$ is computable by binary exponentiation: recursively
compute $R(x) = x^{\lfloor K/2 \rfloor} \bmod \Gamma(x)$, square it modulo
$\Gamma(x)$, and multiply by $x$ once more (modulo $\Gamma(x)$) if $K$ is
odd. Reducing an arbitrary product modulo $\Gamma(x)$ by schoolbook
polynomial remainder costs $O(d^2)$; with fast (FFT-based) polynomial
multiplication and division it drops to $O(M(d))$. Either way, the
squaring step happens $O(\log K)$ times.

\begin{notebox}
The ``multiply by $x$'' step deserves special attention: unlike a general
product modulo $\Gamma(x)$, multiplying a reduced polynomial by the single
monomial $x$ only ever produces one term of overflow (the old
$x^{d-1}$ coefficient becomes an $x^d$ term), so substituting
$x^d \equiv c_1x^{d-1}+\cdots+c_d$ back in costs just $O(d)$ --- far
cheaper than a full $O(M(d))$ multiplication. Structuring the doubling
recursion to isolate this cheap case (rather than folding it into a
generic square-and-multiply loop that treats both branches as full
multiplications) measurably improves the constant factor.
\end{notebox}

\section{Bostan--Mori's Algorithm}

\subsection{Recap: recurrences as rational generating functions}
Writing $A(x) = \sum_i A_i x^i$, the companion lecture shows
$A(x) = P(x)/Q(x)$ with
\[
  Q(x) = 1 - c_1 x - c_2 x^2 - \cdots - c_d x^d, \qquad
  P(x) = A(x)Q(x) \bmod x^d.
\]
(Note $Q$'s coefficients are the same $c_1, \ldots, c_d$ as $\Gamma(x)$
above, but in \emph{reverse} order.) The target is
$A_K = [x^K]\,P(x)/Q(x)$.

\subsection{Halving the target index}
\begin{theorembox}
For any formal power series $Q(x)$, the product $Q(x)Q(-x)$ has zero odd
part: $Q(x)Q(-x) = V(x^2)$ for some power series $V$.

\emph{Why.} For odd $k$, the coefficient $[x^k]Q(x)Q(-x)$ pairs up terms
$Q_a Q_b(-1)^b$ and $Q_b Q_a(-1)^a$ for each even/odd split $a+b=k$
($a$ even, $b$ odd); since $(-1)^b=-1$ and $(-1)^a=1$, these two
contributions are negatives of each other and cancel. \qed
\end{theorembox}

Apply this with $Q(-x)$: multiplying numerator and denominator by
$Q(-x)$,
\[
  \frac{P(x)}{Q(x)} = \frac{P(x)Q(-x)}{Q(x)Q(-x)} = \frac{U(x)}{V(x^2)},
  \qquad U(x) = P(x)Q(-x).
\]
Splitting $U(x) = U_e(x^2) + xU_o(x^2)$ into even/odd parts gives
\[
  [x^K]\frac{P(x)}{Q(x)} =
  \begin{cases}
    [x^{K/2}]\, U_e(x)/V(x) & K \text{ even} \\
    [x^{(K-1)/2}]\, U_o(x)/V(x) & K \text{ odd,}
  \end{cases}
\]
where $U_e, U_o$ have degree $< d$ and $V$ has degree $\le d$ with
$V(0)=1$ --- exactly the same shape of problem, with the target index cut
in (about) half. Iterating until $K=0$ reduces the answer to
$[x^0]P(x)/Q(x) = [x^0]P(x)$, since $Q(0)=1$.

% FIGURE: kth_halving
\begin{figure}[H]
\centering
\begin{tikzpicture}[x=1cm,y=1cm]
\node[lnode,minimum width=1.8cm] (k0) at (0.00,0) {$K=13$\\[-1pt]\scriptsize$(1101)_2$};
\node[lnode,minimum width=1.8cm] (k1) at (3.40,0) {$K=6$\\[-1pt]\scriptsize$(110)_2$};
\node[lnode,minimum width=1.8cm] (k2) at (6.80,0) {$K=3$\\[-1pt]\scriptsize$(11)_2$};
\node[lnode,minimum width=1.8cm] (k3) at (10.20,0) {$K=1$\\[-1pt]\scriptsize$(1)_2$};
\node[lnode,minimum width=1.8cm] (k4) at (13.60,0) {$K=0$\\[-1pt]\scriptsize$(0)_2$};
\draw[larrow] (k0) -- (k1) node[midway,above=1pt,llabel,align=center] {keep the\\ odd part};
\draw[larrow] (k1) -- (k2) node[midway,above=1pt,llabel,align=center] {keep the\\ even part};
\draw[larrow] (k2) -- (k3) node[midway,above=1pt,llabel,align=center] {keep the\\ odd part};
\draw[larrow] (k3) -- (k4) node[midway,above=1pt,llabel,align=center] {keep the\\ odd part};
\node[lnode,lhi,minimum width=3.2cm] (ans) at (13.60,-1.7) {$A_K=[x^0]P(x)=P_0$};
\draw[larrow] (k4) -- (ans);
\node[llabel,align=left,anchor=west] at (0,-1.7) {each round: $P\leftarrow$ parity part of $P(x)Q(-x)$,\\[-1pt] $Q\leftarrow$ even part of $Q(x)Q(-x)$, $K\leftarrow\lfloor K/2\rfloor$};
\end{tikzpicture}
\caption{Bostan--Mori for $K=13=1101_2$. The parity of $K$ at each round decides whether the odd or even part of the numerator is kept; after $\lfloor\log_2 K\rfloor+1=4$ rounds $K$ reaches $0$ and the answer is the constant term of the current numerator.}\label{fig:kth_halving}
\end{figure}

\begin{definitionbox}[(Bostan--Mori algorithm)]
While $K>0$: compute $U(x)=P(x)Q(-x)$ and $W(x)=Q(x)Q(-x)$; replace $P(x)$
with $U$'s even or odd part according to $K$'s parity, replace $Q(x)$
with $W$'s even part, and halve $K$. Once $K=0$, return $[x^0]P(x)$.
\end{definitionbox}

Each round costs two polynomial multiplications of size $O(d)$, and there
are $O(\log K)$ rounds, giving $O(M(d)\log K)$ overall --- the same
asymptotic bound as Fiduccia, though with a smaller constant: standard
estimates put Fiduccia around $3M(d)\log K$ ring operations against
Bostan--Mori's $2M(d)\log K$ (and, with FFT-specific optimizations,
Bostan--Mori's leading term can be pushed down further still). A
follow-up lecture on advanced Bostan--Mori topics covers these constant-
factor refinements.

\section{Reference Implementations}
Both examples below implement Fiduccia's algorithm and the Bostan--Mori
algorithm (with naive $O(d^2)$ polynomial arithmetic, giving
$O(d^2\log K)$ overall) and cross-check them against each other and
against direct recurrence unrolling.

\codenote
\subsection{Rust}
\corelisting{Rust}{example.rs}

\subsection{Python}
\corelisting{Python}{example.py}

\section{Worked Example}
\begin{examplebox}
For the Fibonacci recurrence ($A_0=0, A_1=1, A_i=A_{i-1}+A_{i-2}$), the
generating function is $A(x) = x/(1-x-x^2)$: $Q(x)=1-x-x^2$,
$P(x)=x$. Both reference implementations recover every Fibonacci number
this way for small $K$ (checked against direct unrolling), and --- since
Fiduccia and Bostan--Mori have no way to communicate except by being
correct --- they still agree with each other on $A_{10^{12}} \bmod
998244353$, computed via $O(\log K)$ rounds each, in a fraction of a
second.
\end{examplebox}

\section{Further Reading}
\begin{itemize}[leftmargin=1.5em]
  \item AtCoder Algorithm Lectures index (source of this curriculum):
        \url{https://info.atcoder.jp/entry/algorithm_lectures/index}
  \item C. M. Fiduccia, ``An Efficient Formula for Linear Recurrences'',
        SIAM Journal on Computing 14(1), 1985, pp.\ 106--112.
  \item A. Bostan, R. Mori, ``A Simple and Fast Algorithm for Computing
        the N-th Term of a Linearly Recurrent Sequence'', SOSA 2021,
        pp.\ 118--132.
\end{itemize}

\end{document}
